\podnaslov{Pogojne pričakovane vrednosti}

\begin{definicija}
  \pojem{Pogojna porazdelitev} diskretne slučajne spremenljivke $X$ glede na
  dogodek $B$ je dana z verjetnostjo $P(X=x \such B)$.
  \pojem{Pogojna pričakovana vrednost} slučajne spremenljivke $X$ glede na
  dogodek $B$ je dana z
  \[
	E(X \such B) = \sum_x x P(X = x \such B).
  \]
\end{definicija}

Alternativno lahko izračunamo
\[
  E(X \cdot I_B) = \sum_x x P(X \cdot I_B = x)
  = P(B) \sum_x x P(X = x \such B) = P(B) E(X \such B),
\]
torej $E(X \such B) = E(X \cdot I_B) / P(B)$.
Iz te formule sledi linearnost pogojne pričakovane vrednosti.

\vprasanje{Definiraj pogojno pričakovano vrednost in izpelji alternativno
  izražavo.}

\begin{izrek}
  Naj bo $\{H_1, H_2, \ldots\}$ particija.
  Naj bo $X$ diskretna slučajna spremenljivka.
  Velja
  \[
	E(X) = \sum_k E(X \such H_k) P(H_k).
  \]
\end{izrek}

\begin{proof}
  Računamo
  \[
	\sum_k E(X \such H_k) P(H_k)
	= \sum_k \sum_x x P(X = x \such H_k) P(H_k)
	= \sum_x P(X = x)
	= E(X).
  \]
\end{proof}

\vprasanje{Dokaži formulo za popolno pričakovano vrednost.}

\begin{definicija}
  Naj bosta $\sv{X}, \sv{Y}$ slučajna vektorja z gostoto
  $f_{\sv{X},\sv{Y}}(\sv{x},\sv{y})$.
  Za $\sv{x}$, za katere je $f_{\sv{X}}(\sv{x}) > 0$, definiramo \pojem{pogojno
	gostoto} $\sv{Y}$ glede na $\{\sv{X} = \sv{x}\}$ kot
  \[
	f_{\sv{Y} \such \sv{X} = \sv{x}} (\sv{y})
	= \frac{f_{\sv{X},\sv{Y}}(\sv{x},\sv{y})}{f_{\sv{X}}(\sv{x})}.
  \]
\end{definicija}

\vprasanje{Definiraj pogojno gostoto.}

\begin{izrek}
  Naj bosta $\sv{X}, \sv{Y}$ slučajna vektorja.
  Za $p = \dim \sv{X}$ velja
  \[
	E(g(\sv{Y})) = \int_{\R^p} E(g(\sv{Y}) \such \sv{X} = \sv{x})
	f_{\sv{X}}(\sv{x}) d\sv{x}.
  \]
\end{izrek}

\begin{proof}
  Računamo
  \begin{align*}
	\int_{\R^p} E(g(\sv{Y}) \such \sv{X} = \sv{x}) f_{\sv{X}}(\sv{x}) d\sv{x}
	&= \int_{\R^p} f_{\sv{X}}(\sv{x}) d\sv{x} \int_{\R^q} g(\sv{y}) f_{\sv{Y} \such \sv{X} = \sv{x}}(\sv{y}) d\sv{y} \\
	&= \int_{\R^{p+q}} g(\sv{y}) f_{\sv{X}}(\sv{x}) f_{\sv{Y} \such \sv{X} = \sv{x}}(\sv{y}) d\sv{x} d\sv{y} \\
	&= \int_{\R^{p+q}} g(\sv{y}) f_{\sv{X},\sv{Y}}(\sv{x},\sv{y}) d\sv{x} d\sv{y} \\
	&= E(g(\sv{Y})).
  \end{align*}
\end{proof}

\vprasanje{Dokaži zvezno verzijo formule za popolno pričakovano vrednost.}